\documentclass[11pt]{article}
\usepackage{amsmath,amsthm,amssymb,booktabs}
\usepackage[margin=1in]{geometry}
\newtheorem{theorem}{Theorem}\newtheorem{lemma}[theorem]{Lemma}\newtheorem{proposition}[theorem]{Proposition}
\newtheorem{corollary}[theorem]{Corollary}
\theoremstyle{remark}\newtheorem{remark}[theorem]{Remark}
\newcommand{\Z}{\mathbb{Z}}
\title{Q597: (W$_B$), (LEN) and (W$_a$) are proved. (E-struct) holds, and so (G3$'$) holds under (S)}
\author{Clio}
\date{PROVE 2026-11-27 c1}
\begin{document}\maketitle
\begin{abstract}\small
In case~b, the step of $E$ after the exit, the pair $(E_{k+3},E^L_{k+3})$, the whole back window, and (LEN) $|\mathrm{mid}_E|=|\mathrm{mid}_L|+1$ are proved. With (Merge), (Merge$'$) and (W$_F$) (proved 11-26 c2), this completes (E-struct) in case~b. The same four tools then prove case~a: $J_E$ steps at $R^E_2$, the pair law is the same, the journeys resync, and here the stale set has no local family at all. They derive (W$_a$) with exactly its 8 classes. Six of these are the size-5 classes. The other two are the classes that first appeared at size~6, so the local derivation predicts the census's extension. Lemma~red of the G3$'$ tex then gives \textbf{(G3$'$) for every diverging $B/C$ pair under (S)}.
Four findings.
(1) \textbf{(Q2) is not needed, and the sentence in the brief was right to be doubted.} The pair lemma needs only that the tiles removed by $E_{k+3}$ are present on $E^L$'s board. They are $M_0$ tiles, and the Prop.~rerun template proves their presence. \emph{But not all of them are present:} exactly one triangle of $G(E_{k+3})$ is removed by $E^L_{k+2}$, in every word and every case-b class. Yesterday's coupling table compared full sets and was right only by luck. The corrected lemma uses $G(E_{k+3})\setminus G(E^L_{k+2})$ together with $C(E^L_{k+2})$, and it still separates the census pairs from the bad ones. Its lesion, which drops $E_{k+3}$'s tiles, separates nothing.
(2) \textbf{$J_E$ does not stop at $R^E_{k+3}$}, although that rhombus has $T_E$'s packed shape. The proof is by absence: $L^E_{k+1}=L_{k+2}-\alpha$ (Theorem~bsall) removes a square overlapping $R^E_{k+3}$, and a stopped rhombus would stay on $L^E$'s board.
(3) \textbf{(W$_B$).} The back window built from the proved local facts alone is \emph{exactly} one of the four census representatives, for each $T_L$ and each pair (4/4, exact tile sets).
(4) \textbf{(LEN)}, by Lemma~suffix of the role tex. After step $k+3$ both journeys stand on the same rhombus. Consider the $M_0$ tiles on $J_E$'s board but not on $E^L$'s. Those coming from $J_L$'s prefix or suffix are never touched by $J_E$, by minimality and Lemma~fin. Those removed by $J_E$'s own earlier steps are gone. The rest fall into three local families. Two are $-\alpha$ shadows of tiles $J_E$ removed on class-0 rhombi, killed by Theorem~AW, because $-\alpha=-\delta_{T_L}$ is an ascending edge vector for $T_E$ as well. The third family lies at the contact, and its tiles are already past $T_E$'s removal box at $R_2$, by the monotone offset of Lemma~window.
\end{abstract}

\section{Setting and facts used}
Notation is that of the Mergeprime tex (\texttt{2026-11-26c2-Q597-Mergeprime.tex}): the pair is in case~b, $k\ge0$, $R_2,\dots,R_{k+2}$ is the corridor, $\alpha=\delta_{T_L}$, $a_B=a1$, $a_C=a0$, and $m=k+2$ is the index of the exit, where $E_m$ and $L_m$ are $(0,b)$ on the same hexagon (Prop.~shadow). $E^L$ is $E$'s re-run on $M_0\cdot J_L$, $L^E$ is $L$'s re-run on $M_0\cdot J_E$. Write $R^E_{k+3}$ and $R^L_{k+3}$ for the rhombi created by $E_{k+2}$ and $L_{k+2}$, and $R'$ for the one created by $E^L_{k+2}$. $G(s)$ is the set of non-rhombus tiles removed by a step $s$, $C(s)$ the non-rhombus tiles it creates. $X_j$, $Y_j$ are the boards of $J_E$ and $E^L$ before step $j$, and $p=p_E$.

All of the following are proved.
Lemma~word (each step is a translate of the good configuration of its class and type);
Lemma~fin (Share$_2$ tex: a journey never removes a non-rhombus tile it created; every removed non-rhombus tile is an initial tile, removed once);
Lemma~nr (A1 tex: no re-creation);
Lemma~dinv (first-contact tex);
Lemma~p2 (A1 tex: before contact, the clean journey's rhombi avoid every rhombus of the other journey);
Prop.~shadow, Prop.~rerun, Thm~corr (Share$_2$ tex);
Theorem~bsall (BS$(L)$): $L^E_i=L_{i+1}-\alpha$ for $0\le i\le k+1$, $L^E_{k+2}=L_0+\gamma$, and $J_L$, $L^E$ agree outside their middles;
Theorem~mp, Lemma~persist, Lemma~nostop and Prop.~wb(1) (Mergeprime tex): $E^L_i=E_i-\alpha$ for $2\le i\le k+1$, and $E^L_{k+2}=(0,a_{T_L})$ at $R_{k+2}-\alpha$;
Theorem~AW and Lemma~window (G1c tex);
Lemma~suffix and Lemma~term (role tex, 11-22 c2);
Lemma~win and Lemma~red (G3$'$ tex).
Hypotheses (S) include (Mos), so every board in sight is a Purbhoo mosaic and a tiling.

The prefix fact used in Prop.~rerun(1) is used again here. \emph{$J_L$'s prefix (steps before $p_L$) removes no tile of $W_E$, by minimality of $p_L$, and symmetrically $J_E$'s prefix removes no tile of $W_L$.}

\section{$J_E$ steps after the exit}
\begin{lemma}[shape; finite]\label{lem:shape}
$R^E_{k+3}$ has class~1 for $T_E$. It has the packed shape of nest $T_E$ and of no other nest. $R^L_{k+3}$ has the packed shape of nest $T_L$ only.
\end{lemma}
\begin{proof}
\texttt{s1w1.py}, both $T_L$. This is translation invariant, so one canonical exit rhombus per $T_L$ suffices.
\end{proof}
So the shape alone does not exclude a stop. The engine stops iff \texttt{packed\_at} returns a free, ideal-addable cell, and that is a global condition.

\begin{proposition}[no stop at $R^E_{k+3}$]\label{prop:nostop}
$J_E$ takes a step at $R^E_{k+3}$.
\end{proposition}
\begin{proof}
Suppose $J_E$ stops there. Then $R^E_{k+3}$ is a tile of $M_0\cdot J_E$, which is $L^E$'s initial board.
\emph{$L^E$ does not remove it before step $p_L+k+1$.} A step removes exactly one rhombus, the journey's own current one (the candidate hexagons contain no other rhombus). Before step $p_L$, $L^E$'s rhombi are those of $J_L$'s prefix and its contact rhombus, and by Lemma~p2 (applied with $J_L$ as the clean journey and $J_E$ as the other) none of them is a rhombus of $J_E$, its final rhombus included. From step $p_L$ to $p_L+k+1$, $L^E$ stands on $R^L_0$ (covered by Lemma~p2, since $k\le\mathrm{fc}$ includes the contact) and then on $R_i-\alpha$, $2\le i\le k+2$ (Theorem~bsall: $L^E_{i-1}=L_i-\alpha$). These have class~0, and $R^E_{k+3}$ has class~1, so their $\Lambda_0$ edges differ.
\emph{The contradiction.} $L^E_{k+1}=L_{k+2}-\alpha$ is $T_L$'s $(0,b)$ at $R_{k+2}-\alpha$. Its removed tiles lie on $L^E$'s board (Lemma~word). \texttt{nostop.py} finds, for both $T_L$, a removed square of it overlapping $R^E_{k+3}$ without being equal to it. The board is a tiling, so this is impossible.
\end{proof}
\noindent\texttt{nostop.py} also finds that $L^E_{k+2}$ ($T_L$'s $(1,a_{T_L})$ at $R^L_{k+3}-\alpha$) removes a triangle overlapping $R^E_{k+3}$. That is a second, independent witness. Neither step's removed tiles overlap $C(E_{k+2})$, so the collision really is with the rhombus.

\section{The pair $(E_{k+3},E^L_{k+3})$}
\begin{lemma}[window; finite]\label{lem:w1}
For every case-b contact class and every corridor word $t_2\cdots t_{k+1}\in\{a0,a1\}^k$, and each of the three types $x$ of $E_{k+3}$ at $R^E_{k+3}$, a tile of $G(E_{k+3})$ is never in $G(L_0)\cup G(L_1)\cup\bigcup_{2\le j\le k+1}G(E^L_j)$. It lies in $G(E^L_{k+2})$ for exactly one triangle when $x\in\{b,a_{T_L}\}$, and never when $x=a_{T_E}$.
\end{lemma}
\begin{proof}
Let $\varphi=(-1,-2,1,-2)$, which rises by $1$ at each class-0 a-step (Mergeprime tex, Lemma~persist). Relative to $\varphi(\min R_{k+2})$, the vertices of $G(E_{k+3})$ have $\varphi\in[-6,4]$. Relative to $\varphi(\min R_2)$, those of $G(L_0)\cup G(L_1)$ have $\varphi\in[-6,2]$. Those of $G(E^L_j)$, relative to $\varphi(\min R_j)$, lie in $[-4,1]$ shifted by $|\varphi(\alpha)|\le2$. Equal tiles share vertices, so a coincidence needs $k\le8$ (for $L_0,L_1$) or $k+2-j\le9$ (for $E^L_j$). The geometry of $E^L_j$ against $E_{k+3}$ depends only on the subword $t_j\cdots t_{k+1}$, and every subword of length $\le10$ occurs as a whole word. \texttt{w1.py} checks every word with $k\le10$ over all contact classes: 8188 words. The only hits are the stated triangle in $G(E^L_{k+2})$, for $x\in\{b,a_{T_L}\}$, at every $k$.
\end{proof}

\begin{proposition}[(Q2$''$): presence]\label{prop:q2}
Every tile of $P_1:=G(E_{k+3})\setminus G(E^L_{k+2})$ lies on $Y_{p+k+3}$, $E^L$'s board before step $k+3$.
\end{proposition}
\begin{proof}
Let $h\in P_1$. By Lemma~fin, $h\in M_0$ and $J_E$ removes $h$ only at step $p+k+3$. So $h\in X_p$. By Lemma~dinv, $X_p\setminus Y_p=M_0\setminus M_0\cdot J_L=W_L$, so $h\in Y_p$ unless $J_L$ removes $h$.
\emph{$J_L$ never removes $h$.} Its prefix does not meet $W_E\ni h$. $L_0,L_1$ do not remove $h$ (Lemma~\ref{lem:w1}). $L_2,\dots,L_{k+2}$ remove $H(E_2),\dots,H(E_{k+2})$ (Prop.~shadow), and $J_E$ removes $h$ only at step $k+3$. After its middle, $J_L$'s steps are $L^E$'s (Theorem~bsall). $L^E$ runs on $M_0\cdot J_E$, where $h$ is absent: $J_E$ removed it and never re-creates it (Lemma~nr). So $L^E$ could remove $h$ only after creating it, which Lemma~fin forbids.
\emph{$E^L$ does not remove $h$ at steps $0,\dots,k+1$.} $E^L$'s prefix is $J_E$'s, and that removes $h$ only at $k+3$. The tiles removed by $E^L_0$ and $E^L_1$ lie in $K$ of Prop.~rerun. $K$ consists of non-rhombus tiles created by $J_L$, which are not $M_0$ tiles (a created tile equal to a present one would lie in the removed hexagon, and \texttt{recreate.out} shows that no configuration removes a tile it creates; equality with a removed one is re-creation, excluded by Lemma~nr), together with tiles of $N_E=G(E_0)$, which $J_E$ removes only at step~$0$. For $2\le j\le k+1$, $E^L_j=E_j-\alpha$ (Theorem~mp), and Lemma~\ref{lem:w1} applies. At step $k+2$, $h\notin G(E^L_{k+2})$ by the definition of $P_1$.
\end{proof}

\begin{theorem}[the pair]\label{thm:pair}
$(E_{k+3},E^L_{k+3})$ is $((1,b),(0,a_{T_E}))$ or $((1,a_{T_L}),(0,b))$. Both steps create the same rhombus $R^*$.
\end{theorem}
\begin{proof}
$E_{k+3}$ exists (Prop.~\ref{prop:nostop}). It is one of the three class-1 configurations of $T_E$ at $R^E_{k+3}$, and $E^L_{k+3}$ is $(0,a_{T_E})$ or $(0,b)$ at $R'$ (Prop.~wb(1)). On the tiling $Y_{p+k+3}$ the following are present: $P_1$ (Prop.~\ref{prop:q2}); $P_2=C(E^L_{k+2})$, created at the previous step; and $P_3=G(E^L_{k+3})$ (Lemma~word). So $P_1\cup P_2\cup P_3$ is pairwise compatible. \texttt{pair2.py} goes through the six combinations. The set is compatible exactly for the two stated pairs, and the other four contain an overlapping pair. It also checks that each stated pair creates the same rhombus (\texttt{resync.py}).
\end{proof}
\noindent\emph{Lesion.} If $P_1$ is dropped, all six combinations are compatible (\texttt{pair2.py}, column ``lesion''). So the presence of $E_{k+3}$'s $M_0$ tiles is what selects the pair. Yesterday's \texttt{wb.py} used all of $G(E_{k+3})$, including the triangle $E^L_{k+2}$ had already removed. Its conclusion is the same, but its premise is false.

\section{(W$_B$)}
\begin{theorem}[(W$_B$)]\label{thm:wb}
The back window $(E_{k+2},E_{k+3},E^L_{k+2},E^L_{k+3},L_{k+2},L^E_{k+1},L^E_{k+2})$, translated so that $E_{k+2}$'s least point is $0$, is one of the four representatives of the G3$'$ tex. Hence (Bk) holds (Lemma~win).
\end{theorem}
\begin{proof}
Every entry is now determined by $T_L$, the pair (Theorem~\ref{thm:pair}), and the exit rhombus $R_{k+2}$. $E_{k+2}$ and $L_{k+2}$ are $T_E$'s and $T_L$'s $(0,b)$ at $R_{k+2}$ (Prop.~shadow). $E^L_{k+2}$ is $(0,a_{T_L})$ at $R_{k+2}-\alpha$ (Prop.~wb(1)). $L^E_{k+1}=L_{k+2}-\alpha$. $L^E_{k+2}=L_0+\gamma$ is a translate of $L_0$, a $(1,a_{T_L})$, taken where $L^E_{k+1}$ leaves its rhombus, at $R^L_{k+3}-\alpha$. By Lemma~word it is $T_L$'s $(1,a_{T_L})$ configuration there. \texttt{wbclass.py} builds these seven steps from the local configuration table, for both $T_L$ and both pairs. It compares the exact tile sets (removed and created, rhombi included) with \texttt{windows.pkl}. All 4 agree with the 4 census representatives, one each.
\end{proof}

\section{(LEN)}
\begin{lemma}[the stale set]\label{lem:stale}
Let $k^*=p+k+4$, and let $D=X_{k^*}\setminus Y_{k^*}$. No step of $J_E$ at or after $k^*$ removes a tile of $D$.
\end{lemma}
\begin{proof}
A tile of $D$ is a rhombus, or a non-rhombus tile created by $J_E$, or an $M_0$ tile.
\emph{Rhombi.} A step removes only the journey's own rhombus. At $k^*$ both journeys stand on $R^*$ (Theorem~\ref{thm:pair}), and $J_E$'s later rhombi are created by $J_E$ itself, so none of them is a tile of $D$.
\emph{Created by $J_E$.} By Lemma~fin, $J_E$ never removes them.
\emph{$M_0$ tiles.} Let $g\in M_0\cap D$. Then $J_E$ has not removed $g$ before $k^*$, so $g\in X_p$. Either $g\notin Y_p$, which means $g\in W_L$ (Lemma~dinv), or $g$ was removed by one of $E^L$'s middle steps $E^L_0,\dots,E^L_{k+3}$. (The prefix steps are shared, and $J_E$ has not removed $g$.)
(a) \emph{$g\in W_L$.} Go through $J_L$ as in Prop.~\ref{prop:q2}. If its prefix removes $g$, then $g\notin W_E$, so $J_E$ never removes $g$. If $L_i$ removes $g$ for $2\le i\le k+2$, then $J_E$ removes $g$ at $E_i$, before $k^*$, and only then. If a suffix step (an $L^E$ step) removes $g$, then $g$ is on $L^E$'s board. $L^E$ did not create it, by Lemma~fin, so $g\in M_0\cdot J_E$. Hence $g\notin W_E$, by Lemma~nr. That leaves $g\in G(L_0)\cup G(L_1)$. If also $g\in G(E_0)\cup G(E_1)$, $J_E$ removes it before $k^*$. Otherwise, by \texttt{r1.py}, the 4-offset $f=\mathrm{om}(g,R_2)$ for $T_E$ is already outside the removal box $Q_{T_E}$ in a coordinate that can only move away from it: a $\psi$-coordinate above the maximum, or an $A$-coordinate below the minimum. This holds for every contact class and the $E_1$ that creates $R_2$. By Lemma~window (G1c tex), $f(i)=\mathrm{om}(g,R_i)$ has non-decreasing $\psi$-part and non-increasing $A$-part along any terminating journey, and a removal at step $i$ needs $f(i)\in Q_{T_E}$. So $J_E$ never removes $g$ after step $p+2$. ($Q_{T_E}$ rebuilt locally equals the engine-derived one, \texttt{windowcheck.out}, 5 offsets each.)
(b) \emph{$g\in G(E^L_0)\cup G(E^L_1)$.} Then $g\in K$, and $g$ is an $M_0$ tile of $K$, so $g\in N_E=G(E_0)$. $J_E$ removes it at step $p$.
(c) \emph{$g\in G(E^L_j)$, $2\le j\le k+1$.} $E^L_j=E_j-\alpha$, so $g=u-\alpha$ with $u\in G(E_j)$ removed by $J_E$ at step $p+j$. $R_j$ has class~0, so its edge vectors are $\pm(0,1;0,0)$ and $\pm(0,0;0,1)$, and these include $-\alpha$. Also $-\alpha$ is \emph{ascending for $T_E$}: $\psi_{T_E}$ changes by $+1$ in the relevant coordinate on every removed tile of every configuration (computed with \texttt{psilib}; for $T_E=C$, $v=(0,0;0,-1)$; for $T_E=B$, $v=(0,1;0,0)$). This is the computation of Lemma~pref of the role tex, transposed to the other target. By Theorem~AW, $J_E$ removes $u-\alpha$ at no step other than $p+j$, and $u-\alpha\notin H_{p+j}$. So $J_E$ never removes $g$.
(d) \emph{$g\in G(E^L_{k+2})\cup G(E^L_{k+3})$.} By \texttt{r3.py}, over all contact classes, all words with $k\le6$ and both pairs, these tiles are as follows. Some lie in $C(L_{k+2})$ and are not $M_0$ tiles. Some lie in $G(E_{k+2})\cup G(E_{k+3})$ and are removed by $J_E$ before $k^*$. Exactly one triangle is left, and it has $g+\alpha\in G(E_{k+2})$. $R_{k+2}$ has class~0, so (c)'s argument applies with $j=k+2$. The local geometry of steps $k+2,k+3$ relative to $R_{k+2}$ does not depend on the word (the window starts at $R_{k+2}$), so $k\le6$ covers every $k$. (A tile of $G(E^L_{k+3})$ in $C(E^L_{k+2})$ is excluded by Lemma~fin for $E^L$.)
\end{proof}

\begin{theorem}[(LEN) for $E$]\label{thm:len}
$|\mathrm{mid}_E|=|\mathrm{mid}_{E^L}|=k+4=|\mathrm{mid}_L|+1$.
\end{theorem}
\begin{proof}
Apply Lemma~suffix (role tex) to $A=J_E$ on $X_j$ and $A'=E^L$ on $Y_j$, both toward $T_E$ on Purbhoo mosaics, at step $k^*$, where both stand on $R^*$. Hypothesis (i) is Lemma~\ref{lem:stale}. (T) is Lemma~term: $J_E$ is a journey toward $T_E$ from a nest-$A$ corner, $E^L$ is its re-run on the board after $J_L$, toward $T_L\ne T_E$, and every rhombus of a journey has the nest-$A$ class or $T_E$'s class~1. So $E^L$ agrees with $J_E$ from $k^*$ on, step for step, and both stop at the same step. The steps $E_{k+3}$ and $E^L_{k+3}$ differ (different rhombi: class~1 against class~0). $E_0\ne E^L_0$ (Prop.~rerun(2): different types). So the middle of $E$, between the longest common prefix and the longest common suffix, is exactly steps $p,\dots,p+k+3$ on both sides. $|\mathrm{mid}_L|=k+3$ is (M2) (Prop.~shadow(3)).
\end{proof}

\section{Case a}\label{sec:casea}
In case~a, Theorem~share2 with Theorem~h1 and Prop.~rerun (Share$_2$ tex) give $E_1=(0,b)$, $L_0=(1,a_{T_L})$, $L_1=(1,a_{T_E})$ and $E^L_1=(1,b)$, and $E^L_0$, all determined by the contact class (4 classes, Cor.~wf). Theorem~bsall with $C=(L_1)$ gives $L$'s middle $L_0L_1$, and $L^E$'s middle $(L_1-\alpha)(L_0+\gamma)$. So $E_1$ plays the part of the exit, and $E^L_1$ the part of $E^L_{k+2}$. Let $R^E_2$ and $R^{EL}_2$ be the rhombi created by $E_1$ and $E^L_1$. \texttt{casea.py} runs the four checks of the case-b proof on the four classes.

\begin{theorem}[case a]\label{thm:casea}
Assume (S) and case~a.
\begin{enumerate}
\item $R^E_2$ has class~1 for $T_E$ and $R^{EL}_2$ has class~0. $J_E$ steps at $R^E_2$, and $E^L$ steps at $R^{EL}_2$.
\item $(E_2,E^L_2)$ is $((1,b),(0,a_{T_E}))$ or $((1,a_{T_L}),(0,b))$, and the two steps create the same rhombus.
\item $|\mathrm{mid}_E|=|\mathrm{mid}_{E^L}|=3$.
\item The four middles form a translate of one of the 8 (W$_a$) representatives. Hence (A) holds and $\Delta_E=\Delta_L=0$ (Lemma~win, Lemma~red(1)).
\end{enumerate}
\end{theorem}
\begin{proof}
(1) For $E^L$: class~0, Lemma~nostop. For $J_E$: suppose it stops at $R^E_2$. Then that rhombus lies on $L^E$'s initial board $M_0\cdot J_E$. Before step $p_L$, $L^E$'s rhombi are $J_L$'s prefix rhombi and the contact rhombus $R^L_0$, and none of them is a rhombus of $J_E$ (Lemma~p2). At $p_L$, $L^E_0=L_1-\alpha$ is taken at $R^L_0$. (\texttt{casea.py} asserts $R^L_1-\alpha=R^L_0$, and that $L^E_0$ removes exactly $G(L_1)-\alpha$.) It removes a square overlapping $R^E_2$, which is impossible on a tiling. $L^E_1$ removes an overlapping triangle too, a second witness.
(2) Let $h\in G(E_2)\setminus(G(L_0)\cup G(L_1))$. As in Prop.~\ref{prop:q2}: $h\in M_0$, and $J_E$ removes it only at $p+2$, so $h\in X_p$. $J_L$ never removes $h$: its prefix does not meet $W_E$, $L_0$ and $L_1$ do not remove $h$ by choice of $h$, and its suffix consists of $L^E$'s steps on $M_0\cdot J_E$, where $h$ is absent. So $h\in Y_p$. $E^L_0$ and $E^L_1$ remove only tiles of $K$ and of $K'=(K\setminus G(E^L_0))\cup C(E^L_0)$ (Prop.~rerun). The $M_0$ tiles of these sets lie in $N_E=G(E_0)$, and $h\notin G(E_0)$. Tiles of $C(E^L_0)$ are created on a board where $h$ is present outside the removed hexagon, so none of them is $h$. Hence $h$ is on $Y_{p+2}$. Also on $Y_{p+2}$ are $K''=(K'\setminus G(E^L_1))\cup C(E^L_1)$ and $G(E^L_2)$. \texttt{casea.py} tests all $3\times3$ type combinations. Exactly the two stated pairs are compatible, in every class. Both create the same rhombus.
(3) Lemma~suffix at $k^*=p+3$, with (T) from Lemma~term. For (i), run through $D\cap M_0$ as in Lemma~\ref{lem:stale}. Prefix tiles of $J_L$: $J_E$ never removes them. Suffix tiles of $J_L$: these are not in $W_E$. The middle of $L$ is $L_0L_1$, and \texttt{casea.py} finds every $M_0$ tile of $G(L_0)\cup G(L_1)$ removed by $E_0$, $E_1$ or $E_2$, in every class and for both pairs. $E^L_0$ and $E^L_1$ remove $M_0$ tiles only from $G(E_0)$. Every $M_0$ tile of $G(E^L_2)$ lies in $G(E_0)\cup G(E_1)\cup G(E_2)$. (The tiles of $G(E^L_2)$ created by $L_0,L_1,E^L_0,E^L_1$ are not $M_0$ tiles: Lemma~nr together with the tiling argument.) So the $M_0$ part of $D$ has no local family, and (i) holds. Rhombi, and tiles created by $J_E$, are handled as in case~b. The steps at $p+2$ differ (class~1 against class~0 rhombi), and so do those at $p$.
(4) \texttt{casea.py} builds the four middles from (1)--(3) and the local configuration table. It compares them, tile for tile, with \texttt{windows.pkl} in the shared-square frame. 4 classes and 2 pairs give 8 middles: 6 equal the six size-5 representatives, and 2 are new to \texttt{windows.pkl}. The two new ones have typed words $E$: \texttt{1b 0b 1a1}, $E^L$: \texttt{1a1 1b 0b} ($T_L=B$) and $E$: \texttt{1b 0b 1a0}, $E^L$: \texttt{1a0 1b 0b} ($T_L=C$). These are exactly the two classes that first appeared on size-6 chunks (G3$'$ tex, \texttt{win6new.out}), where identity (A) was checked. \texttt{casea.py} matches only against the size-5 pickle. The two size-6 classes are identified by their typed words, not by tile sets. That is a weaker match, recorded in the grading below.
\end{proof}

\begin{corollary}[(E-struct), and (G3$'$)]\label{cor:g3}
Assume (S). Every diverging $B/C$ pair satisfies (E-struct): (W$_a$) in case~a, and (LEN), (Merge), (Merge$'$), (W$_F$), (W$_B$) in case~b. Hence $\Delta_E=\Delta_L$ for every diverging $B/C$ pair (Lemma~red of the G3$'$ tex), which is (G3$'$).
\end{corollary}
\begin{proof}
Case~a: Theorem~\ref{thm:casea}(4). Case~b: (LEN) is Theorem~\ref{thm:len}. (Merge) on the corridor is Prop.~shadow(2) of the Share$_2$ tex. (Merge$'$) is Theorem~mp. (W$_F$) is Cor.~wf. (W$_B$) is Theorem~\ref{thm:wb}.
\end{proof}

\section{Grading the chain (read link by link)}
Each input of Lemma~red was re-read against the registry \texttt{q585-ybe-train-vs-commutor.json} before this file called anything proved.
\begin{center}\small\begin{tabular}{lll}\toprule
input of Lemma red & source & registry trust\\\midrule
Lemma delta, Lemma win (16 reps) & G3$'$ tex & proved (\texttt{Q597-G3prime-reduction} children)\\
(LEN), $m'_E=m_E$, case b & Thm~\ref{thm:len} here & new: \texttt{Q597-LEN-E}\\
$m'_L=m_L=k+3$ & Thm bsall, Prop.~shadow(3) & proved (\texttt{Q597-Merge-M2})\\
(Merge), $E_i=L_i$, $2\le i\le k+1$ & Prop.~shadow(2) & proved (\texttt{Q597-Merge-corr})\\
(Merge$'$) & Thm mp & proved (\texttt{Q597-Mergeprime})\\
(W$_F$) & Cor.~wf & proved (\texttt{Q597-WF})\\
(W$_B$) & Thm~\ref{thm:wb} here & new: \texttt{Q597-WB-class}; \texttt{Q597-WB-pair} upgraded\\
(W$_a$) & Thm~\ref{thm:casea} here & new: \texttt{Q597-Wa-struct}\\\bottomrule
\end{tabular}\end{center}
Links inside today's proofs that rest on earlier files: Lemma~fin, Prop.~rerun, Thm~h1 (Share$_2$ tex, under \texttt{Q597-Share2}); Thm~bsall; Lemma~p2 and Lemma~nr (A1 tex); Thm~AW and Lemma~window (G1c tex); Lemma~suffix and Lemma~term (role tex); Lemma~dinv (first-contact tex). The 11-27 WAKE brief lists all of these as proved. \textbf{Two points are cited in a form slightly different from their statements:}
(a) Lemma~p2 is stated for ``the clean journey'' and ``$J_o$''. I apply it with $J_L$ as the clean journey and $J_E$ as $J_o$. Its proof uses only the minimality of the first contact, (S1) and Lemma~nr, so it holds with the roles as used. This is YELLOW, not RED.
(b) Lemma~term says ``in particular (T) holds''. That needs every rhombus of $J_E$ to have the nest-$A$ class or $T_E$'s class~1. This is Lemma~class, inside registry node \texttt{Q597-both-diverge} (proved), checked today against the registry. CLOSED.
(c) The two size-6 (W$_a$) classes are matched by typed word only. So I do not lean on that match. \texttt{identities.py} evaluates (A) \emph{directly} on all 8 locally built case-a middles, and (Bk) on the 4 local back windows. It uses $\delta$ from \texttt{lookstep.py}, with no census involved: $\Delta_E=\Delta_L=0$ on 8/8, and (Bk) holds on 4/4. Control: translating one step by a unit vector makes the identity fail in 18 of 24 cases. In the other 6 the translation moves no readout key, which is Lemma~delta(2). So case~a and the back window rest on Lemma~delta and this direct evaluation, not on the representatives table. CLOSED.
(S) itself, including (Mos), is a standing hypothesis. (Mos) is computed, not proved.

\section{Computation}
Code is in \texttt{proofs/code-1127-Q597/}. The prediction was written before any run (\texttt{PREDICTION.md}).
\paragraph{Prediction against outcome.} W1: predicted to hold with the window closing early. \textbf{Partly wrong:} one triangle of $G(E_{k+3})$ is removed by $E^L_{k+2}$ in every word, so the natural (Q2$'$) (``all of $G(E_{k+3})$ is present'') is false. S1: I gave 60\% that $R^E_{k+3}$ is \emph{not} of packed shape. \textbf{Wrong}: it is $T_E$'s packed shape, and the no-stop needed the $L^E$ absence argument.
\paragraph{Local, finite (all translation invariant, canonical frames or the case-b/case-a contact classes of \texttt{contacts()}).}
\begin{center}\small\begin{tabular}{lll}\toprule
script & statement & result\\\midrule
\texttt{s1w1.py} & shapes of $R^E_{k+3}$, $R^L_{k+3}$ & packed for $T_E$ resp.\ $T_L$ only\\
\texttt{nostop.py} & $L^E_{k+1}$, $L^E_{k+2}$ overlap $R^E_{k+3}$ & a square, resp.\ a triangle, both $T_L$\\
\texttt{w1.py} & Lemma~\ref{lem:w1}, $k\le10$ & 8188 words, hits only $E^L_{k+2}$ (one triangle)\\
\texttt{pair2.py} & pair lemma on $Y_{p+k+3}$ & exactly the 2 census pairs; lesion (no $P_1$): 6/6 survive\\
\texttt{pair2.py} lesion 2 & uncorrected premise, all of $G(E_{k+3})$ & \textbf{0/6} survive (premise inconsistent)\\
\texttt{resync.py} & same rhombus after step $k+3$ & 4/4\\
\texttt{wbclass.py} & local back windows vs \texttt{windows.pkl} & 4/4 exact\\
\texttt{r1.py} & $G(L_0)\cup G(L_1)$ past $Q_{T_E}$ at $R_2$ & 4 tiles past, 20 removed by $E_0,E_1$\\
\texttt{r3.py} & leftover tile of $E^L_{k+2},E^L_{k+3}$ & one triangle, $g+\alpha\in G(E_{k+2})$, $k\le6$\\
\texttt{psilib} one-liner & $-\alpha$ ascending for $T_E$ & $\Delta\psi=+1$ on every removed tile, both\\
\texttt{casea.py} & case a, all four checks & 4 classes, 8 middles = 8 (W$_a$) classes\\
\texttt{identities.py} & (A) and (Bk) on the local windows & 8/8 and 4/4; control fails 18/24\\\bottomrule
\end{tabular}\end{center}
\paragraph{Census, size~5, all 24 chunks} (\texttt{stale\_e.py}, \texttt{stale5.out}; case~b, $k=0$ included, 521 records). (LEN) 521/521. (Q2) as stated yesterday (tiles of $E^L_{k+3}$ outside $C(L_{k+2})$ are $M_0$ tiles of $G(E_{k+3})$): 1042/1042. (Q2$''$): in every record exactly one tile of $G(E_{k+3})$ is removed by $E^L_{k+2}$, two by $E^L_{k+3}$, and none by $J_L$. The (W$_B$) pairs at $k=0$ are the same four classes (41, 91, 41, 92). This answers Step~0(a). The stale set $D$ at $k^*$ falls only into the families of Lemma~\ref{lem:stale}: $J_L$-prefix tiles (rhombi included), $L_0$ triangles, $J_L$-suffix tiles, $E^L$ corridor tiles, and the $E^L_{k+2}$ triangle. \textbf{None is ever removed by $J_E$ later}, 0 exceptions.
\paragraph{Out of sample, size~6} (fresh chunks 13--15 of 24, never used before). \texttt{stale6.out}: 584 case-b records (198 with $k=0$, 386 with $k>0$, $k$ up to $6$). (LEN) 584/584. (Q2) 1168/1168. (Q2$''$): one tile by $E^L_{k+2}$ and two by $E^L_{k+3}$ in every record, none by $J_L$. Back-window pairs: the same four classes (26, 68, 41, 63 at $k=0$; 69, 135, 60, 122 at $k>0$). Stale set: 10\,191 tiles, \textbf{0} ever removed by $J_E$ later.
\paragraph{$B/B$ control.} Statement predicted to fail on $B/B$: (LEN), and with it the exit divergence on which Prop.~\ref{prop:nostop} rests. For $T_E=T_L$ the exit hexagon is rotated the same way by both journeys, so $L^E_{k+1}$ has nothing to collide with. Known: (LEN) 0/440 on $B/B$ (G3$'$ tex) and (M3) 249/249 shared steps after the exit. Not re-run today. The proofs here use $T_E\ne T_L$ in Lemma~term and in Lemma~\ref{lem:shape}.

\section{Status}
\begin{itemize}
\item \textbf{Proved today (under (S)):} Prop.~\ref{prop:nostop} (no stop after the exit), Prop.~\ref{prop:q2} (Q2$''$), Theorem~\ref{thm:pair} (the pair), Theorem~\ref{thm:wb} (W$_B$), Lemma~\ref{lem:stale}, Theorem~\ref{thm:len} ((LEN) for $E$), Theorem~\ref{thm:casea} (case a, (W$_a$)), Corollary~\ref{cor:g3} ((E-struct), hence (G3$'$)).
\item \textbf{Not needed:} (Q2) in yesterday's form. It is true (1042/1042), but the proof routes around it. Yesterday's sentence (``a provenance statement \dots which the absence form cannot avoid here'') was half right. Provenance was needed, but of $E_{k+3}$'s tiles, not $E^L_{k+3}$'s, and it is the Prop.~rerun template, a presence statement. Absence then finishes it.
\item \textbf{YELLOW link} (grading (a)): Lemma~p2 applied with the roles exchanged. Its proof is role-symmetric.
\item \textbf{Standing hypothesis:} (S), with (Mos) computed only.
\item \textbf{Next:} (G3$'$) $\Rightarrow$ (MC) for the pair, given PQ-RIGID and $\mathcal F$-membership (G3$'$ tex corollary). Re-read that corollary's inputs before promoting anything above (G3$'$).
\end{itemize}
\end{document}
